\RequirePackage[2023-06-01]{latexrelease}
\documentclass[a4paper,11pt,noarxiv]{quantumarticle}
\usepackage{amsmath,amssymb,amsthm}
\usepackage{graphicx}
\usepackage{booktabs}
\usepackage{array}
\usepackage{enumitem}
\usepackage{xcolor}
\definecolor{linkink}{HTML}{1c3f77}
\usepackage{hyperref}
\usepackage{microtype}
\emergencystretch=1.5em
\hypersetup{colorlinks=true,linkcolor=linkink,citecolor=linkink,
urlcolor=linkink}

\newcommand{\R}{\mathbb{R}}
\newcommand{\Rm}{R^{\mathrm m}}
\newcommand{\Risk}{\operatorname{Risk}}
\newcommand{\norm}[1]{\lVert #1\rVert}
\newcommand{\abs}[1]{\lvert #1\rvert}
\newcommand{\dd}{\,\mathrm{d}}
\newcommand{\grad}{\nabla}
\newcommand{\fbar}{\bar f}
\DeclareMathOperator{\capa}{cap}
\DeclareMathOperator{\dist}{dist}
\DeclareMathOperator{\Ric}{Ric}
\newcommand{\status}[1]{\textsc{\small [#1]}}

\theoremstyle{plain}
\newtheorem{theorem}{Theorem}
\newtheorem{proposition}[theorem]{Proposition}
\newtheorem{lemma}[theorem]{Lemma}
\newtheorem{corollary}[theorem]{Corollary}
\theoremstyle{definition}
\newtheorem{example}[theorem]{Example}
\theoremstyle{remark}
\newtheorem{remark}[theorem]{Remark}

\title{Shrinkage priors, admissibility, and weighted Riemannian
geometry: potential theory and co-dimension two obstructions}
\author{Oleg Roshka}
\affiliation{Independent Researcher \textnormal{--}
\href{https://orcid.org/0009-0009-3075-9417}{ORCID
0009-0009-3075-9417}}
\date{24 October 2026}

\begin{document}
\maketitle

\begin{abstract}
Two literatures intersect without citing each other: Stein-type
shrinkage and admissibility in statistical decision theory (Stein,
Brown, Komaki) and the potential theory of weighted Riemannian
manifolds (Grigor'yan; Perelman's weighted scalar curvature;
Saloff-Coste). This note establishes the exact dictionary between
them and resolves the analytical question that their intersection
raises. First, on a statistical manifold $(M,g)$ equipped with the
Fisher--Rao metric, Komaki's $O(N^{-2})$ Kullback--Leibler risk
reduction of a prior $e^{-\Phi}\pi_J$ over the Jeffreys prior $\pi_J$
equals $(\Rm_\Phi-R)/(2N^2)$, the $\Phi$-dependent part of Perelman's
weighted scalar curvature $\Rm_\Phi=R+2\Delta_g\Phi-\abs{\grad\Phi}_g^2$,
while the prior-independent $O(N^{-2})$ risk depends on statistical
embedding curvatures and skewness rather than the intrinsic scalar
curvature $R$. Second, Brown's (1971) admissibility characterization
for normal means $X\sim N_m(\theta,I)$ under a generalized Bayes
marginal $f^*=p*F$ is the parabolic potential theory of the weighted
Euclidean manifold $W_F=(\R^m,\delta_{ij},f^*\dd x)$, and his 1971
sphere-average test is the sphere-area parabolicity test attributed in
Riemannian geometry to Lyons--Sullivan (1984) and Grigor'yan (1985).
Third, we investigate when the doubling and scale-invariant Poincar\'e
hypotheses of heat-kernel theory allow Grigor'yan's ball-volume test
to decide admissibility problems on which both of Brown's radial tests
are silent. We prove that a bounded spherical $A_2$ ratio
$Q(r)=\fbar(r)A(r)$ makes Brown's two radial tests exhaustive, and
that in one dimension doubling and the Poincar\'e inequality force this
bound, so the ball-volume test decides nothing new for product weights.
In dimension $m\ge3$, however, this implication fails: because a
one-dimensional curve in $\R^3$ has co-dimension two, it does not
disconnect balls, allowing $Q(r)\to\infty$ while preserving uniform
doubling and Poincar\'e constants. Exploiting this co-dimension two
phenomenon, we construct an explicit doubling and Poincar\'e weight on
$\R^3$ on which both of Brown's radial tests fail, yet Grigor'yan's
ball-volume criterion decides inadmissibility. Finally, for a
borderline non-doubling product weight left undecided by all three
geometric tests, we prove inadmissibility analytically via Krein's
spectral string theory and Kasahara's Tauberian theorem.
\end{abstract}

\section{Introduction and setting}\label{sec:intro}

In decision-theoretic statistics, the phenomenon discovered by Stein
\cite{stein1956}---that the standard maximum likelihood estimator for a
$d$-dimensional normal mean is inadmissible for $d\ge3$---is governed
by the recurrence or transience of an associated diffusion operator.
For quadratic loss in the normal location model $X\sim N_m(\theta,I)$,
Brown \cite{brown1971} proved that a generalized Bayes estimator with
bounded drift is admissible if and only if a symmetric diffusion on
$\R^m$ weighted by the marginal density $f^*=p*F$ is recurrent. For
Kullback--Leibler predictive loss on a general curved statistical
manifold $(M,g)$, Komaki \cite{komaki2006} proved that a Bayesian
predictive density built from a prior $f=e^{-\Phi}\pi_J$ asymptotically
dominates the Jeffreys predictive density if and only if
$u=\sqrt{f/\pi_J}=e^{-\Phi/2}$ is nonconstant positive superharmonic
with respect to the Laplace--Beltrami operator $\Delta_g$ of the
Fisher--Rao metric.

Independently, in geometric analysis and potential theory on smooth
metric measure spaces $(M,g,e^{-\Phi}\dd V)$, the same differential
operators and integral criteria were developed to study heat kernels,
parabolicity (recurrence of weighted Brownian motion), and Ricci flow
with a potential \cite{bakryemery1985,grigoryan1999,perelman2002}.
Despite identical analytical objects, the statistical and geometric
literatures developed largely in parallel.

The purpose of this note is twofold:
\begin{enumerate}[label=\textup{(\roman*)},leftmargin=2em,itemsep=2pt]
\item To record the precise dictionary between the two fields
(\S\S\ref{sec:komaki}--\ref{sec:brown} and Table~\ref{tab:dict}),
identifying both exact identities and structural boundaries where naive
analogies fail.
\item To resolve the analytical question opened by this dictionary
(\S\ref{sec:reach}): does the modern Grigor'yan--Saloff-Coste--Sturm
theory of doubling metric measure spaces supporting a scale-invariant
Poincar\'e inequality yield an admissibility criterion (the ball-volume
test) that strictly extends Brown's 1971 radial sphere-average and
ray tests?
\end{enumerate}

We show that the answer to (ii) depends on a topological co-dimension
threshold. In one dimension---and consequently for one-dimensional
product priors on $\R^m$---the doubling and Poincar\'e hypotheses force
the arithmetic and harmonic sphere means of $f^*$ to remain comparable
(Lemmas~\ref{lem:exhaustive}--\ref{lem:1d_a2} and
Corollary~\ref{cor:product}), which renders Brown's two radial tests
already exhaustive. In dimension $m\ge3$, however, zeros or valleys of
the weight along a submanifold of co-dimension at least two do not
disconnect balls. Diffusing paths can bypass co-dimension two
obstructions without inflating the scale-invariant Poincar\'e constant,
even as the spherical $A_2$ ratio diverges
(Theorem~\ref{thm:codim2_unbounded}). By placing co-dimension two rings
across dyadic shells so that their radial projections sweep the unit
sphere infinitely often, we construct a smooth doubling and Poincar\'e
weight on $\R^3$ on which every radial ray integral diverges and the
sphere-average integral converges---leaving both of Brown's tests
silent---while Grigor'yan's ball-volume test immediately establishes
transience (Theorem~\ref{thm:blind_brown}). We conclude in
Proposition~\ref{prop:krein} by settling a borderline non-doubling
product example analytically using Feller's canonical scale and
Kasahara's Karamata--Tauberian theorem for Krein strings.

\paragraph{Notation and standing definitions.}
On the parametric model side, let $\{p(x\mid\theta):\theta\in M\subset\R^d\}$
be a regular statistical model carrying the Fisher--Rao metric
$g_{ij}(\theta)=E_\theta[\partial_i\ell\,\partial_j\ell]$, where
$\ell=\log p(x\mid\theta)$, the Jeffreys prior density
$\pi_J(\theta)=\abs{\det g(\theta)}^{1/2}$ with Riemannian volume
element $\dd V=\pi_J(\theta)\dd\theta$, and the Laplace--Beltrami
operator $\Delta_g=\abs{g}^{-1/2}\partial_i(\abs{g}^{1/2}g^{ij}\partial_j)$
\cite{amari2000}. A smooth positive prior density is written
$f(\theta)=e^{-\Phi(\theta)}\pi_J(\theta)$, with amplitude
$u(\theta)=\sqrt{f(\theta)/\pi_J(\theta)}=e^{-\Phi(\theta)/2}$. The
triple $(M,g,e^{-\Phi}\dd V)$ is a smooth metric measure space with
weighted (Bakry--\'Emery) Laplacian
$\Delta_\Phi=\Delta_g-\grad\Phi\cdot\grad$, self-adjoint on
$L^2(M,e^{-\Phi}\dd V)$, Bakry--\'Emery Ricci tensor
$\Ric_\Phi=\Ric_g+\grad^2\Phi$ \cite{bakryemery1985}, and Perelman's
weighted scalar curvature \cite{perelman2002}
\begin{equation}\label{eq:perelman}
\Rm_\Phi=R+2\Delta_g\Phi-\abs{\grad\Phi}_g^2,
\end{equation}
where $R=g^{ij}\Ric_{ij}$ is the Levi-Civita scalar curvature of
$(M,g)$. Integrating $\Rm_\Phi$ against $e^{-\Phi}\dd V$ yields
Perelman's entropy functional
$\mathcal F(g,\Phi)=\int_M\Rm_\Phi\,e^{-\Phi}\dd V
=\int_M(R+\abs{\grad\Phi}_g^2)e^{-\Phi}\dd V$ whenever the boundary
term of integration by parts vanishes. A weighted manifold is called
\emph{parabolic} if its weighted Brownian motion with generator
$\Delta_\Phi$ is recurrent, or equivalently if it admits no positive
Green function and every positive superharmonic function
($\Delta_\Phi h\le0$) is constant \cite[\S5]{grigoryan1999}. We denote
by $\Risk_N(f)$ the expected Kullback--Leibler divergence from the true
density $p(y\mid\theta)$ to the Bayesian predictive density
$\hat p_f(y\mid x^{(N)})$ formed from $N$ independent observations
under prior $f$.

On Brown's normal estimation side, let $X\sim N_m(\theta,I_m)$ in
$\R^m$ ($m\ge3$), let $F$ be a $\sigma$-finite prior measure on $\R^m$,
and let $f^*(x)=(p*F)(x)=\int_{\R^m}(2\pi)^{-m/2}e^{-\abs{x-\theta}^2/2}\dd F(\theta)$
be its marginal density. Under quadratic loss
$E_\theta[\abs{\delta(X)-\theta}^2]$, the generalized Bayes estimator is
Tweedie's score estimator
$\delta_F(x)=x+\grad\log f^*(x)=x+\gamma_F(x)$, which has bounded risk
whenever $\gamma_F=\grad\log f^*$ is bounded \cite[\S3]{brown1971}.
This defines the weighted Euclidean space
$W_F=(\R^m,\delta_{ij},f^*\dd x)$ with measure $\dd\mu=f^*\dd x$, dual
measure $\dd\nu=(f^*)^{-1}\dd x$, and weighted Laplacian
$\Delta_\Phi=\Delta+\grad\log f^*\cdot\grad$ ($\Phi=-\log f^*$), which
generates Brown's diffusion $\{Z_t\}$. For $r>0$, let $S^{m-1}$ denote
the unit sphere with normalized surface measure $\dd\sigma$, and define
the spherical arithmetic mean, harmonic-inverse mean, and
\emph{spherical $A_2$ ratio}:
\begin{equation}\label{eq:spherical_a2}
\fbar(r)=\int_{S^{m-1}}f^*(r\varphi)\dd\sigma(\varphi),\quad
A(r)=\int_{S^{m-1}}\frac{\dd\sigma(\varphi)}{f^*(r\varphi)},\quad
Q(r)=\fbar(r)A(r)\ge1.
\end{equation}

\section{Predictive shrinkage and Perelman's weighted scalar curvature}
\label{sec:komaki}

We first record the exact relation between Komaki's asymptotic risk
expansion \cite{komaki2006,komaki2015} and Perelman's weighted scalar
curvature \eqref{eq:perelman}.

\begin{proposition}\label{prop:komaki_perelman}
For a regular statistical model $(M,g)$ and a smooth prior
$f=e^{-\Phi}\pi_J$ with $u=e^{-\Phi/2}$, the $O(N^{-2})$ predictive risk
reduction $\varrho_N(f)=\Risk_N(\pi_J)-\Risk_N(f)$ satisfies
\begin{equation}\label{eq:komaki_perelman}
\varrho_N(f)
=\frac{1}{2N^2}\bigl(2\Delta_g\Phi-\abs{\grad\Phi}_g^2\bigr)+o(N^{-2})
=\frac{\Rm_\Phi-R}{2N^2}+o(N^{-2})
=-\frac{2}{N^2}\,\frac{\Delta_g u}{u}+o(N^{-2}).
\end{equation}
Consequently, $f$ asymptotically dominates $\pi_J$ up to $O(N^{-2})$ if
and only if $u$ is a nonconstant positive superharmonic function on
$(M,g)$, or equivalently $\Rm_\Phi\ge R$ everywhere with strict
inequality at some point.
\end{proposition}

\begin{proof}
In Komaki's asymptotic expansion of the predictive risk \cite[eqs.\
(3)--(5)]{komaki2006}, subtracting $\Risk_N(f)$ from $\Risk_N(\pi_J)$
cancels all prior-independent terms and third-order skewness
contractions, leaving the first equality of \eqref{eq:komaki_perelman}
written in terms of $\log(f/\pi_J)=-\Phi$. The second equality is
\eqref{eq:perelman}, and the third follows from the identity
$\Delta_g(e^{-\Phi/2})/e^{-\Phi/2}=-\tfrac12\Delta_g\Phi+\tfrac14\abs{\grad\Phi}_g^2$.
The domination criterion is Komaki's Theorem~2 \cite{komaki2006}.
\end{proof}

\begin{proposition}\label{prop:no_R}
In static independent sampling, the prior-independent $O(N^{-2})$ term
of the Jeffreys predictive risk $\Risk_N(\pi_J)$ is not proportional to
the Riemannian scalar curvature $R$ of $(M,g)$, and Perelman's
functional $\mathcal F(g,\Phi)$ is not the integrated predictive risk.
Whenever the boundary term of integration by parts vanishes, up to
$o(N^{-2})$,
\begin{equation}\label{eq:int_risk}
\int_M\varrho_N(f)\,e^{-\Phi}\dd V
=\frac{1}{2N^2}\int_M\abs{\grad\Phi}_g^2\,e^{-\Phi}\dd V
=\frac{1}{2N^2}\Bigl[\mathcal F(g,\Phi)-\int_M R\,e^{-\Phi}\dd V\Bigr].
\end{equation}
\end{proposition}

\begin{proof}
Komaki \cite[Theorem A1, eq.\ (29)]{komaki2015} derives the complete
$O(N^{-2})$ expansion of $\Risk_N(\pi_J)$. Its prior-independent part is
a linear combination of the exponential and mixture statistical
embedding curvatures, the skewness norm $T_{ijk}T^{ijk}$, a fourth-order
cumulant contraction, and a mixed moment term. For a one-dimensional
curved exponential family ($d=1$), these embedding curvatures and
cumulants are generically nonzero (recovering Efron's statistical
curvature \cite{efron1975}), whereas the intrinsic Riemannian scalar
curvature $R$ vanishes identically for every $d=1$ manifold. Equation
\eqref{eq:int_risk} follows from \eqref{eq:komaki_perelman} and the
integration by parts identity
$\int_M(\Delta_g\Phi)e^{-\Phi}\dd V=\int_M\abs{\grad\Phi}_g^2e^{-\Phi}\dd V$.
\end{proof}

\begin{remark}[Role of curvature and conformal change of vague prior]
\label{rem:conformal}
By Komaki's Theorems~1--2 \cite{komaki2006}, a shrinkage prior
asymptotically dominating $\pi_J$ exists if and only if $(M,g)$ is
non-parabolic (carries a positive Green function). Riemannian curvature
enters static shrinkage domination through this global potential-theoretic
threshold---for instance, negative sectional curvature in $d\ge3$ implies
non-parabolicity \cite[Theorem~3]{komaki2006}---rather than as an
additive potential $R$ in the local risk. Furthermore, when $d\ge3$ and
one tests domination over a general baseline prior $h=e^{-\Phi}\pi_J$
rather than $\pi_J$, Komaki \cite[\S3, Theorems 4--6]{komaki2006} shows
that $f$ dominates $h$ if and only if $\sqrt{f/h}$ is superharmonic
with respect to the conformally rescaled metric
$\tilde g=(h/\pi_J)^{2/(d-2)}g=e^{-2\Phi/(d-2)}g$. Under this conformal
change of metric, standard transformation rules for scalar curvature
yield, up to boundary terms,
\begin{equation}\label{eq:conformal_F}
\mathcal F(g,\Phi)=\int_M\Bigl[\tilde R-\frac{1}{d-2}
\abs{\tilde\grad\Phi}_{\tilde g}^2\Bigr]\dd\tilde V.
\end{equation}
For the flat normal location family $N_d(\mu,I_d)$ and Stein's harmonic
prior $h(\mu)=\norm{\mu}^{-(d-2)}$, the conformal metric $\tilde g_{ij}
=\norm{\mu}^{-2}\delta_{ij}$ is isometric to the flat Riemannian
cylinder $S^{d-1}\times\R$ under $\mu\mapsto(\mu/\norm{\mu},\ln\norm{\mu})$,
pushing the origin $\mu=0$ to infinite geodesic distance
\cite[Example~1]{komaki2006}.
\end{remark}

\section{Brown (1971) as potential theory on \texorpdfstring{$W_F$}{W\_F}}
\label{sec:brown}

In the normal location model $X\sim N_m(\theta,I_m)$, Stein's unbiased
risk identity writes the quadratic risk difference between $\delta_F=x+\gamma_F$
and the maximum likelihood estimator $\delta_0(x)=x$ as
$E_\theta[\abs{\delta_F(X)-\theta}^2]-m=E_\theta[2\operatorname{div}\gamma_F(X)+\abs{\gamma_F(X)}^2]
=4E_\theta[\Delta(f^*)^{1/2}/(f^*)^{1/2}]$. Brown's fundamental
contribution \cite{brown1971} was to characterize exact (non-asymptotic)
admissibility of $\delta_F$ in terms of the Dirichlet form
$\mathcal E_{W_F}(j)=\int_{\R^m}\abs{\grad j}^2f^*\dd x$ on the weighted
manifold $W_F=(\R^m,\delta_{ij},f^*\dd x)$.

\begin{theorem}[Brown \cite{brown1971} in weighted manifold terminology]
\label{thm:brown}\leavevmode
\begin{enumerate}[label=\textup{(\alph*)},leftmargin=1.8em,itemsep=2pt,topsep=2pt]
\item \textup{(Capacity criterion; Brown Thm.\ 4.3.1, Grigor'yan
\cite[Thm.\ 5.1]{grigoryan1999})} The diffusion $\{Z_t\}$ on $W_F$ is
transient if and only if the weighted capacity of the closed unit ball
$B_1=\{\abs{x}\le1\}$,
\[
\capa_{W_F}(B_1)=\inf\Bigl\{\int_{\R^m}\abs{\grad j(x)}^2f^*(x)\dd x :
j\in C_c^\infty(\R^m),\; j\big|_{B_1}=1\Bigr\},
\]
is strictly positive, i.e.\ $W_F$ is non-parabolic.
\item \textup{(Admissibility $\Leftrightarrow$ parabolicity; Brown Thm.\
5.1.1)} If $\gamma_F=\grad\log f^*$ is bounded, then $\delta_F$ is
admissible if and only if $W_F$ is parabolic. Conversely, if $W_F$ is
transient, then $\delta_F$ is inadmissible without any boundedness
assumption on $\gamma_F$.
\item \textup{(Spherically symmetric case; Brown Thm.\ 6.3.1, Grigor'yan
\cite[Prop.\ 3.1]{grigoryan1999})} If $F$ is spherically symmetric,
$\delta_F$ is inadmissible if and only if
$\int_1^\infty r^{1-m}f^*(r)^{-1}\dd r<\infty$.
\item \textup{(Sphere-average parabolicity test; Brown Thm.\ 6.4.3,
Lyons--Sullivan \cite{lyonssullivan1984}, Grigor'yan \cite{grigoryan1985}
\cite[Thm.\ 7.5]{grigoryan1999})} For any prior $F$, if
\begin{equation}\label{eq:sphere_test}
\int_1^\infty\frac{\dd r}{r^{m-1}\fbar(r)}=\infty,
\end{equation}
then $W_F$ is parabolic, and hence $\delta_F$ is admissible whenever
$\gamma_F$ is bounded. Furthermore, by the Cauchy--Schwarz inequality on
dyadic shells, \eqref{eq:sphere_test} is implied by (and is strictly
stronger than) Grigor'yan's ball-volume parabolicity condition
$\int_1^\infty r\dd r/\mu(B_r)=\infty$ \cite[Thm.\ 7.3]{grigoryan1999}.
\item \textup{(Ray transience test; Brown Thm.\ 6.4.4)} If there exists
a measurable subset $\Omega\subset S^{m-1}$ of positive surface measure
such that
\begin{equation}\label{eq:ray_test}
\sup_{\varphi\in\Omega}\int_1^\infty\frac{\dd r}{r^{m-1}f^*(r\varphi)}<\infty,
\end{equation}
then $W_F$ is transient and $\delta_F$ is inadmissible.
\item \textup{(1/$r$ drift equivalence; Brown Thm.\ 6.4.2)} Under the
decay condition $\abs{\gamma_F(x)}\le k/\abs{x}$ for large $\abs{x}$,
condition \eqref{eq:sphere_test} is both necessary and sufficient for
admissibility.
\end{enumerate}
\end{theorem}

\begin{remark}[Historical priority on the sphere-area test]
In the Riemannian geometry literature, the sufficiency of the
reciprocal sphere-area integral \eqref{eq:sphere_test} for parabolicity
is attributed to Lyons and Sullivan (1984) \cite{lyonssullivan1984} and
Grigor'yan (1985) \cite{grigoryan1985}; see \cite[p.~178]{grigoryan1999}.
Brown's Theorem~6.4.3 \cite{brown1971} establishes the identical
variational test on $W_F$ thirteen years earlier by plugging the radial
test function $j(r)=\big(\int_r^R s^{1-m}\fbar(s)^{-1}\dd s\big)\big/
\big(\int_1^R s^{1-m}\fbar(s)^{-1}\dd s\big)$ into the capacity
infimum of Theorem~\ref{thm:brown}(a). Table~\ref{tab:dict} summarizes
the complete dictionary between the statistical and geometric
formulations.
\end{remark}

\begin{table*}[t]
\caption{Dictionary between statistical decision theory / information
geometry and weighted Riemannian geometry.}
\label{tab:dict}
\small
\begin{tabular}{@{}>{\raggedright\arraybackslash}p{0.46\textwidth}>{\raggedright\arraybackslash}p{0.50\textwidth}@{}}
\toprule
Statistical Theory (Stein, Brown, Komaki) & Weighted Riemannian Geometry (Grigor'yan, Perelman) \\
\midrule
Parametric model with Fisher metric $(M,g)$ & Riemannian manifold $(M,g)$ \\
Smooth prior $f=e^{-\Phi}\pi_J$ & Smooth metric measure space $(M,g,e^{-\Phi}\dd V)$ \\
$O(N^{-2})$ KL risk reduction $\varrho_N(f)$ over $\pi_J$ & $(\Rm_\Phi-R)/(2N^2)=-2\Delta_g u/(N^2u)$ (Prop.~\ref{prop:komaki_perelman}) \\
Asymptotic domination over $\pi_J$ (Komaki Thm.\ 2) & $u=e^{-\Phi/2}$ nonconstant positive superharmonic \\
Existence of dominating shrinkage prior (Komaki Thm.\ 1) & $(M,g)$ non-parabolic (carries a Green function) \\
Prior-independent $O(N^{-2})$ risk of $\pi_J$ & Embedding curvatures and skewness; \emph{not} $R$ (Prop.~\ref{prop:no_R}) \\
Change of baseline vague prior $h=e^{-\Phi}\pi_J$ ($d\ge3$) & Conformal metric $\tilde g=e^{-2\Phi/(d-2)}g$ (Remark~\ref{rem:conformal}) \\
Stein prior $\norm{\mu}^{-(d-2)}$ in conformal frame & Flat Riemannian cylinder $S^{d-1}\times\R$ \\
Generalized Bayes admissibility, bounded $\gamma_F$ (Brown Thm.\ 5.1.1) & Parabolicity of $W_F=(\R^m,\delta_{ij},f^*\dd x)$ \\
Brown's variational energy infimum (Brown Thm.\ 4.3.1) & Weighted capacity $\capa_{W_F}(B_1)$ (Grigor'yan Thm.\ 5.1) \\
Brown's sphere-average test (Brown Thm.\ 6.4.3, 1971) & Sphere-area parabolicity test (Grigor'yan Thm.\ 7.5, 1985) \\
--- & Ball-volume parabolicity test $\int r\dd r/\mu(B_r)=\infty$ (weaker) \\
Doubling + Poincar\'e ball-volume transience (Thm.~\ref{thm:blind_brown}) & Decides co-dimension 2 weights where Brown's tests fail \\
\bottomrule
\end{tabular}
\end{table*}

\section{The reach of Brown's radial tests and co-dimension two obstructions}
\label{sec:reach}

Theorem~\ref{thm:brown}(d) is a sufficient condition for parabolicity
(recurrence) and Theorem~\ref{thm:brown}(e) is a sufficient condition
for transience. As Brown observed \cite[p.~900]{brown1971}, there exist
anisotropic priors $F$ for which the sphere-average integral
\eqref{eq:sphere_test} converges while the ray integral in
\eqref{eq:ray_test} diverges along almost every direction
$\varphi\in S^{m-1}$, rendering both tests silent.

In geometric analysis, on the other hand, a celebrated theorem of
Grigor'yan \cite{grigoryan1991} and Saloff-Coste \cite{saloffcoste1992}
(extended to local Dirichlet spaces by Sturm \cite{sturm1996}) states
that if a weighted manifold satisfies the volume doubling property
$\mu(B(x,2\rho))\le D\,\mu(B(x,\rho))$ and the scale-invariant $(2,2)$-Poincar\'e
inequality
\begin{equation}\label{eq:poincare}
\int_{B(x,\rho)}\abs{u-u_B}^2\dd\mu
\le C\rho^2\int_{B(x,\rho)}\abs{\grad u}^2\dd\mu
\end{equation}
uniformly for all balls $B=B(x,\rho)$, then two-sided Gaussian heat
kernel bounds hold, and transience of $W_F$ is \emph{equivalent} to the
convergence of the ball-volume integral:
\begin{equation}\label{eq:vol_transience}
W_F \text{ is transient} \quad\Longleftrightarrow\quad
\int_1^\infty\frac{r\dd r}{\mu(B(0,r))}<\infty
\qquad\text{\cite[Thm.\ 11.1]{grigoryan1999}}.
\end{equation}
A natural question is whether \eqref{eq:vol_transience} resolves
admissibility problems on which Brown's radial tests (d) and (e) are
silent. We now determine the exact conditions under which this occurs.

\subsection{Exhaustiveness under bounded spherical \texorpdfstring{$A_2$}{A2}
and collapse for product weights}

\begin{lemma}[Bounded spherical $A_2$ makes Brown's tests exhaustive]
\label{lem:exhaustive}
Suppose the spherical $A_2$ ratio \eqref{eq:spherical_a2} satisfies
$\sup_{r\ge1}Q(r)=C<\infty$. Then Brown's two radial tests are
exhaustive: if $\int_1^\infty r^{1-m}\fbar(r)^{-1}\dd r=\infty$, then
$W_F$ is parabolic; if $\int_1^\infty r^{1-m}\fbar(r)^{-1}\dd r<\infty$,
then $W_F$ is transient and $\delta_F$ is inadmissible.
\end{lemma}

\begin{proof}
The first case is Theorem~\ref{thm:brown}(d). In the second case, let
$I(\varphi)=\int_1^\infty r^{1-m}f^*(r\varphi)^{-1}\dd r$. By Tonelli's
theorem, its average over $\varphi\in S^{m-1}$ is
\[
\int_{S^{m-1}}I(\varphi)\dd\sigma(\varphi)
=\int_1^\infty\frac{A(r)}{r^{m-1}}\dd r
\le C\int_1^\infty\frac{\dd r}{r^{m-1}\fbar(r)}<\infty.
\]
Hence $I(\varphi)<\infty$ for $\sigma$-almost every $\varphi\in S^{m-1}$.
By Chebyshev's inequality, the set $\Omega=\{\varphi\in S^{m-1}:
I(\varphi)\le 2\int_{S^{m-1}}I\dd\sigma\}$ has surface measure at least
$1/2$, so Theorem~\ref{thm:brown}(e) applies.
\end{proof}

Lemma~\ref{lem:exhaustive} implies that any weight on which both of
Brown's radial tests are silent must satisfy $\limsup_{r\to\infty}Q(r)=\infty$.
In dimension one, however, doubling and the scale-invariant Poincar\'e
inequality forbid such divergence.

\begin{lemma}[One-dimensional doubling and Poincar\'e force $A_2$ on
centred intervals]\label{lem:1d_a2}
Let $w>0$ be continuous on $\R$, $\dd\mu=w(t)\dd t$, and
$\dd\nu=w(t)^{-1}\dd t$. Suppose $\mu$ is doubling on intervals,
$\mu(2I)\le D\,\mu(I)$, and satisfies the one-dimensional
scale-invariant Poincar\'e inequality \eqref{eq:poincare} on all
intervals $I$ with constant $C$. Then for every $r>0$,
\begin{equation}\label{eq:1d_a2}
\mu([-r,r])\,\nu([-r,r])\le 32\,CD^3\,r^2.
\end{equation}
\end{lemma}

\begin{proof}
On $I=[-2r,2r]$ (of length $\abs{I}=4r$), define the continuous piecewise
smooth test function $u(t)=0$ for $t\in[-2r,-r]$, $u(t)=1$ for
$t\in[r,2r]$, and $u(t)=\nu([-r,t])/\nu([-r,r])$ for $t\in[-r,r]$. This
choice minimizes the Dirichlet energy across $[-r,r]$, giving
$\int_I\abs{u'(t)}^2\dd\mu(t)=\nu([-r,r])^{-1}$. On the other hand, any
function equal to $0$ on a set of $\mu$-measure $a=\mu([-2r,-r])$ and
equal to $1$ on a set of $\mu$-measure $b=\mu([r,2r])$ has $\mu$-variance
on $I$ at least $ab/(a+b)\ge\tfrac12\min(a,b)$. Since $[-r,r]$ is
covered by three successive doublings of $[r,2r]$ and of $[-2r,-r]$,
doubling yields $\min(a,b)\ge D^{-3}\mu([-r,r])$. Substituting both
bounds into \eqref{eq:poincare} gives
$\tfrac12 D^{-3}\mu([-r,r])\le C(4r)^2\nu([-r,r])^{-1}$, which is
\eqref{eq:1d_a2}.
\end{proof}

\begin{corollary}[Collapse for one-dimensional product weights]
\label{cor:product}
Let $m\ge3$ and suppose $f^*(x)=w(x_1)$ (or its Gaussian convolution
when $\abs{(\log w)'}$ is bounded), where $w$ satisfies the hypotheses
of Lemma~\ref{lem:1d_a2}. Then $\sup_{r\ge1}Q(r)<\infty$, so Brown's
tests are exhaustive by Lemma~\ref{lem:exhaustive}; in fact, $\delta_F$
is always inadmissible by Theorem~\ref{thm:brown}(e), and the
ball-volume criterion \eqref{eq:vol_transience} decides nothing beyond
Brown's radial tests.
\end{corollary}

\begin{proof}
For $m\ge3$, the spherical means of $w(x_1)$ over $S^{m-1}(r)$ reduce to
weighted integrals over $x_1\in[-r,r]$ with Jacobian kernel
$c_m(1-(x_1/r)^2)^{(m-3)/2}r^{-1}$, which is bounded above on $[-r,r]$
and bounded below on $[-r/\sqrt{2},r/\sqrt{2}]$. By doubling,
$\fbar(r)\asymp\mu([-r,r])/r$ and $A(r)\lesssim\nu([-r,r])/r$, so
$Q(r)\lesssim\mu([-r,r])\nu([-r,r])/r^2\le 32CD^3$ by
Lemma~\ref{lem:1d_a2}. Furthermore, any doubling measure on $\R$
satisfies reverse doubling $\mu([-r,r])\ge c\,r^\varepsilon$ for some
$\varepsilon>0$, whence $\int_1^\infty r^{1-m}\fbar(r)^{-1}\dd r
\lesssim\int_1^\infty r^{2-m-\varepsilon}\dd r<\infty$ for $m\ge3$.
Lemma~\ref{lem:exhaustive} then implies inadmissibility via
Theorem~\ref{thm:brown}(e).
\end{proof}

\subsection{Co-dimension two obstructions in \texorpdfstring{$\R^m$ ($m\ge3$)}{Rm (m>=3)}}

Why does Lemma~\ref{lem:1d_a2} hold in one dimension? Because on $\R$, a
point $t=0$ where $w(t)$ is small has co-dimension one: it disconnects
the interval $[-2r,2r]$ into two halves, forcing every continuous path
from $[-2r,-r]$ to $[r,2r]$ to pass through the bottleneck at $t=0$.
For example, $w(t)=(1+\abs{t})^2$ is doubling on $\R$ but fails the
scale-invariant Poincar\'e inequality because of the bottleneck at $t=0$,
and its interval ratio grows as $r$.

In dimension $m\ge3$, the topology changes fundamentally: a
one-dimensional smooth curve $\Gamma\subset\R^m$ has \emph{co-dimension
$k=m-1\ge2$}. Removing a co-dimension $\ge2$ set does not disconnect
any ball in $\R^m$; test functions and diffusing trajectories can pass
around $\Gamma$ through the bulk where the weight is large. Meanwhile,
where $\Gamma$ pierces the $(m-1)$-dimensional sphere $S^{m-1}(r)$, the
intersection $\Gamma\cap S^{m-1}(r)$ also has co-dimension $m-1\ge2$
inside the sphere, allowing $(f^*)^{-1}$ to blow up rapidly enough to
drive $Q(r)\to\infty$ without violating either doubling or the
scale-invariant Poincar\'e inequality on balls.

\begin{theorem}[Breakdown of Lemma~\ref{lem:1d_a2} in dimension $m\ge3$]
\label{thm:codim2_unbounded}
For every $m\ge3$ and every $\alpha>m-1$, the cylindrical co-dimension
$(m-1)$ weight on $\R^m$ measuring distance to the $x_m$-axis
$\Gamma_0=\{x_1=\dots=x_{m-1}=0\}$,
\begin{equation}\label{eq:cyl_weight}
w(x)=\bigl(1+x_1^2+\dots+x_{m-1}^2\bigr)^{\alpha/2},
\end{equation}
satisfies uniform volume doubling and the scale-invariant $(2,2)$-Poincar\'e
inequality \eqref{eq:poincare} on all Euclidean balls $B(x_0,\rho)\subset\R^m$,
yet its spherical $A_2$ ratio diverges at a polynomial rate:
\begin{equation}\label{eq:q_div}
Q(r)=\fbar(r)A(r)\asymp r^{\alpha-(m-1)}\xrightarrow{r\to\infty}\infty.
\end{equation}
\end{theorem}

\begin{proof}
Write $x=(y,x_m)$ with $y\in\R^{m-1}$. Since $w(x)=(1+\abs{y}^2)^{\alpha/2}$
is independent of $x_m$, doubling and the scale-invariant Poincar\'e
inequality on balls in $\R^m$ follow by tensorization from the
corresponding properties of $\tilde w(y)=(1+\abs{y}^2)^{\alpha/2}$ on
$\R^{m-1}$. At scales $\rho\le1$, $\tilde w$ is smooth and bounded above
and below on any ball of radius $\rho$ by a constant ratio; at scales
$\rho\ge1$, $\tilde w(y)$ is uniformly comparable to the homogeneous
power weight $\abs{y}^\alpha$ outside the unit ball. Since the origin
$y=0$ is a single point of co-dimension $k=m-1\ge2$ in $\R^{m-1}$ and
$\alpha>0>-k$, $\abs{y}^\alpha$ is a $p$-admissible weight in the sense
of Heinonen--Kilpel\"ainen--Martio for all $p\ge1$, satisfying uniform
doubling (with constant $D\le 2^{m+\alpha}$) and the scale-invariant
$(2,2)$-Poincar\'e inequality on all balls in $\R^{m-1}$
\cite{kilpelainen1994,bjorn2011}.

To evaluate $Q(r)$, parameterize $S^{m-1}(r)$ by $\abs{y}=r\sin\theta$
with $\theta\in[0,\pi]$. Then
\[
\fbar(r)=c_m\int_0^{\pi/2}\bigl(1+r^2\sin^2\theta\bigr)^{\alpha/2}
\sin^{m-2}\theta\dd\theta\asymp r^\alpha,
\]
whereas for $\alpha>m-1$ the integral for $A(r)$ localizes to the polar
caps $\theta\lesssim r^{-1}$:
\[
A(r)=c_m\int_0^{\pi/2}\bigl(1+r^2\sin^2\theta\bigr)^{-\alpha/2}
\sin^{m-2}\theta\dd\theta
\asymp\int_0^1(1+r^2\theta^2)^{-\alpha/2}\theta^{m-2}\dd\theta
\asymp r^{-(m-1)}.
\]
Multiplying the two asymptotics yields \eqref{eq:q_div}.
\end{proof}

We now use this co-dimension two mechanism to answer the question posed
at the beginning of \S\ref{sec:reach}: we construct a smooth doubling
and Poincar\'e weight on $\R^3$ where the co-dimension two valley winds
across shells so as to block every radial ray from the origin, silencing
both of Brown's radial tests while Grigor'yan's ball-volume criterion
\eqref{eq:vol_transience} decides transience.

\begin{theorem}[A doubling--Poincar\'e weight decided by ball volume
where both of Brown's radial tests fail]\label{thm:blind_brown}
First, for any parabolic weighted manifold $W_F$ satisfying doubling and
the scale-invariant Poincar\'e inequality, Grigor'yan's volume test
$\int_1^\infty r\dd r/\mu(B_r)=\infty$ implies $\int_1^\infty r^{1-m}\fbar(r)^{-1}\dd r=\infty$
by Theorem~\ref{thm:brown}(d), so Brown's sphere-average test is never
silent in the parabolic regime.

In the transient regime on $\R^3$, however, there exist smooth weights
$f^*>0$ with bounded logarithmic gradient $\abs{\grad\log f^*}\le C$ such
that:
\begin{enumerate}[label=\textup{(\roman*)},leftmargin=2em,itemsep=2pt]
\item $W_F=(\R^3,\delta_{ij},f^*\dd x)$ satisfies uniform volume
doubling and the scale-invariant Poincar\'e inequality \eqref{eq:poincare}
on all balls $B(x_0,\rho)\subset\R^3$;
\item Brown's sphere-average test \eqref{eq:sphere_test} is silent:
$\int_1^\infty r^{-2}\fbar(r)^{-1}\dd r<\infty$;
\item Brown's ray test \eqref{eq:ray_test} is silent along \emph{every}
direction: $\int_1^\infty r^{-2}f^*(r\varphi)^{-1}\dd r=\infty$ for all
$\varphi\in S^2$;
\item Grigor'yan's ball-volume test \eqref{eq:vol_transience} holds,
$\int_1^\infty r\dd r/\mu(B(0,r))<\infty$, proving that $W_F$ is
transient and $\delta_F$ is inadmissible.
\end{enumerate}
\end{theorem}

\begin{proof}
The parabolic statement follows immediately from the Cauchy--Schwarz
inequality on dyadic shells $B_{2^k}\setminus B_{2^{k-1}}$:
\[
\Bigl(2^{k-1}\Bigr)^2
=\Bigl(\int_{2^{k-1}}^{2^k}\dd r\Bigr)^2
\le\Bigl(\int_{2^{k-1}}^{2^k}\omega_{m-1}r^{m-1}\fbar(r)\dd r\Bigr)
\Bigl(\int_{2^{k-1}}^{2^k}\frac{\dd r}{\omega_{m-1}r^{m-1}\fbar(r)}\Bigr),
\]
so $\int_{2^{k-1}}^{2^k}r^{1-m}\fbar(r)^{-1}\dd r\ge c\,2^{2k}/\mu(B_{2^k})
\asymp\int_{2^{k-1}}^{2^k}r\dd r/\mu(B_r)$.

To construct the transient example on $\R^3$, fix $\alpha\ge2$. For each
integer $k\ge4$, divide the dyadic shell $r\in[2^k,2^{k+1}]$ into
$N_k=2^k$ sub-shells of radius $r_{k,j}=2^k+j$ ($j=0,1,\dots,2^k-1$). On
each sphere $S^2(r_{k,j})$, let $\Gamma_{k,j}\subset S^2(r_{k,j})$ be the
parallel of latitude $\theta_{k,j}=\pi(j+\tfrac12)/2^k$ in spherical
coordinates $(r,\theta,\phi)$ (for $j$ near $0$ or $2^k$, a small circle
of geodesic radius $\pi/2^k$ around the pole). Then each $\Gamma_{k,j}$
is a smooth closed one-dimensional circle in $\R^3$ (hence of
co-dimension two in $\R^3$), and distinct circles in the collection
$\Gamma=\bigcup_{k\ge4}\bigcup_{j=0}^{2^k-1}\Gamma_{k,j}$ are mutually
separated in three-dimensional Euclidean distance by at least $1$, while
successive dyadic blocks $k$ and $k'$ can be separated by leaving buffer
shells $r\in[\tfrac32 2^k,2^{k+1}]$ empty (reducing $N_k$ to $2^{k-1}$).
Now define the smooth weight
\begin{equation}\label{eq:ring_weight}
f^*(x)=\biggl(\frac{1+\dist(x,\Gamma)^2}{1+\abs{x}^2}\biggr)^{\alpha/2}
\qquad\text{on }\R^3,
\end{equation}
smoothed on the unit scale near the cutoffs of the distance function so
that $\abs{\grad\log f^*(x)}\lesssim(1+\dist(x,\Gamma))^{-1}\le1$ is
globally bounded.

We verify properties (i)--(iv) in turn:
\begin{itemize}[leftmargin=1.5em,itemsep=2pt]
\item \emph{Property (ii) and (iv):} On the empty buffer shells
$r\in[\tfrac32 2^k,2^{k+1}]$, $\dist(x,\Gamma)\asymp\abs{x}$, so
$f^*(x)\asymp1$. Even within the active shells, $\dist(x,\Gamma)\le1$
only in a tubular neighborhood of total volume $O(2^k\cdot 2^k)=O(r^2)$
inside a shell of volume $\asymp r^3$, and elsewhere $\dist(x,\Gamma)$
grows linearly away from the rings. Consequently, $\mu(B(0,r))\asymp r^3$
and $\fbar(r)\asymp1$, which yields both
$\int_1^\infty r^{-2}\fbar(r)^{-1}\dd r<\infty$ and
$\int_1^\infty r\dd r/\mu(B(0,r))<\infty$.
\item \emph{Property (iii):} Let $\varphi=(\theta,\phi)\in S^2$ be any
direction. For every $k\ge4$, there exists an index $j\in\{0,\dots,N_k-1\}$
such that $\abs{\theta-\theta_{k,j}}\le\pi/2^k$. As the ray $r\varphi$
crosses the radius $r_{k,j}\asymp 2^k$, its Euclidean distance to the
circle $\Gamma_{k,j}$ satisfies $\dist(r\varphi,\Gamma_{k,j})\le\abs{r-r_{k,j}}
+r_{k,j}\abs{\theta-\theta_{k,j}}\le 1+2\pi$ for all
$r\in[r_{k,j}-1,r_{k,j}+1]$. On this radial interval of length $2$,
$f^*(r\varphi)\asymp(1+\abs{x}^2)^{-\alpha/2}\asymp 2^{-k\alpha}$. Thus,
the ray integral across scale $k$ contributes at least
\[
\int_{r_{k,j}-1}^{r_{k,j}+1}\frac{\dd r}{r^2 f^*(r\varphi)}
\gtrsim 2\cdot 2^{-2k}\cdot 2^{k\alpha}=2\cdot 2^{k(\alpha-2)}\ge 2
\qquad(\text{since }\alpha\ge2).
\]
Summing over $k\ge4$ forces $\int_1^\infty r^{-2}f^*(r\varphi)^{-1}\dd r=\infty$
for \emph{every} $\varphi\in S^2$.
\item \emph{Property (i):} Let $B=B(x_0,\rho)$ be any ball in $\R^3$ at
scale $\abs{x_0}\asymp 2^k$. Because the circles $\Gamma_{k,j}$ are
$1$-separated and have radius of curvature $\gtrsim 2^k$ (or $\ge1$ near
the poles), if $\rho\le\tfrac12$, $B$ intersects the tubular
neighborhood of at most one curve $\Gamma_{k,j}$. On $B$, the denominator
$(1+\abs{x}^2)^{\alpha/2}\asymp 2^{k\alpha}$ is constant up to a factor
of $2^\alpha$, which cancels out of both the doubling ratio
$\mu(2B)/\mu(B)$ and the scale-invariant Poincar\'e inequality
\eqref{eq:poincare}. In local tubular coordinates around the
co-dimension two curve $\Gamma_{k,j}$, the numerator
$(1+\dist(x,\Gamma_{k,j})^2)^{\alpha/2}$ is bi-Lipschitz equivalent to
the cylindrical co-dimension two model weight $(1+y_1^2+y_2^2)^{\alpha/2}$
of Theorem~\ref{thm:codim2_unbounded}, which satisfies uniform doubling
and scale-invariant Poincar\'e on all balls. Finally, for larger radii
$\rho>\tfrac12$, the co-dimension two tubes where $f^*\ll1$ occupy a
small fraction of the cross-sectional area and never disconnect $B$, so
standard John-domain chaining across the bulk where
$f^*(x)\asymp(\rho/2^k)^\alpha$ preserves uniform doubling and
scale-invariant Poincar\'e constants.
\end{itemize}
By \eqref{eq:vol_transience} and Theorem~\ref{thm:brown}(b), $W_F$ is
transient and $\delta_F$ is inadmissible.
\end{proof}

\subsection{Analytic Krein--Tauberian resolution of a borderline
non-doubling product weight}

When a weight fails the doubling condition as well, all three geometric
tests---Theorem~\ref{thm:brown}(d), Theorem~\ref{thm:brown}(e), and
Grigor'yan's volume test \eqref{eq:vol_transience}---can fail
simultaneously. For product priors, we can settle the borderline case
analytically via Krein's spectral string theory \cite{kasahara1975,kotani1982}.

\begin{proposition}[Analytic resolution of a logarithmic product prior]
\label{prop:krein}
On $\R^3$, let $F(\dd\theta)=(1+\abs{\theta_1})^{-1}\log(e+\abs{\theta_1})
\dd\theta_1\dd\theta_2\dd\theta_3$, so that $f^*(x)=w(x_1)$ with
$w(t)\asymp(1+\abs{t})^{-1}\log(e+\abs{t})$ and bounded drift
$\abs{\gamma_F}\le C$. Then:
\begin{enumerate}[label=\textup{(\roman*)},leftmargin=2em,itemsep=2pt]
\item Theorem~\ref{thm:brown}(d) is silent ($\fbar(r)\asymp r^{-1}\log^2 r$,
so $\int_1^\infty r^{-2}\fbar(r)^{-1}\dd r<\infty$), Theorem~\ref{thm:brown}(e)
is silent ($\int_1^\infty r^{-2}w(r\cos\theta)^{-1}\dd r=\infty$ for all
$\cos\theta\ne0$), and $w$ is not doubling ($\mu([0,2r])/\mu([r/2,3r/2])\asymp\log r\to\infty$).
\item The one-dimensional diffusion $W_t$ on $\R$ with generator
$\mathcal L_W u=w(x)^{-1}(w(x)u'(x))'$ is null-recurrent, and its
transition density $p_t^W(0,0)$ with respect to the invariant measure
$w(x)\dd x$ satisfies the exact asymptotic
\begin{equation}\label{eq:krein_asymp}
p_t^W(0,0)\sim\frac{C_0}{(\log t)^2}\qquad\text{as }t\to\infty
\qquad(C_0>0).
\end{equation}
\item Consequently, the three-dimensional product diffusion
$Z_t=(W_{2t},B_{2t}^{(2)})$ on $W_F$ is transient, and $\delta_F$ is
inadmissible.
\end{enumerate}
\end{proposition}

\begin{proof}
Part (i) follows by direct integration: $\mu([0,r])=\int_0^r w(t)\dd t
\asymp\tfrac12(\log r)^2$, while on any interval $[r/2,3r/2]$ away from
the origin $\mu([r/2,3r/2])\asymp\log r$, so the doubling ratio grows as
$\log r$.

For parts (ii) and (iii), write the one-dimensional generator
$\mathcal L_W=\frac{\dd}{\dd M}\frac{\dd}{\dd s}$ in Feller's canonical
scale $s(x)$ and speed measure $M(x)$ on $[0,\infty)$ (by symmetry,
reflecting at $0$ scales the return density by $2$):
\[
s(x)=\int_0^x\frac{\dd y}{w(y)}
\asymp\int_0^x\frac{(1+y)\dd y}{\log(e+y)}
\sim\frac{x^2}{2\log x}\qquad(x\to\infty).
\]
Since $s(\infty)=\infty$, $W_t$ is recurrent on $\R$. Inverting $s(x)$
gives $x(s)\sim\sqrt{s\log s}$ as $s\to\infty$, so
$\log x(s)=\tfrac12\log s+O(\log\log s)$. Expressing the speed measure
$M$ as a function of the natural scale $s$ yields
\[
\widetilde M(s)=M(x(s))=\int_0^{x(s)}w(y)\dd y
\asymp\frac{1}{2}\bigl(\log x(s)\bigr)^2
\sim\frac{1}{8}(\log s)^2\qquad(s\to\infty).
\]
As a function of natural scale $s$, $\widetilde M(s)\sim\tfrac18(\log s)^2$
is \emph{slowly varying at infinity} in Karamata's sense (regularly
varying of index $\alpha=0$: $\lim_{s\to\infty}\widetilde M(\lambda s)/
\widetilde M(s)=1$ for all $\lambda>0$). By Kasahara's Tauberian theorem
for Krein strings in natural scale \cite[Theorem~2]{kasahara1975} (see
also Kotani--Watanabe \cite{kotani1982}), when $\widetilde M(s)$ is
slowly varying at infinity, the Laplace transform of the return density
$G_\lambda(0,0)=\int_0^\infty e^{-\lambda t}p_t^W(0,0)\dd t$ satisfies
$G_\lambda(0,0)\sim(2\lambda\widetilde M(1/\lambda))^{-1}$ as
$\lambda\to0^+$, and Karamata's Tauberian theorem yields
$p_t^W(0,0)\sim(2\widetilde M(t))^{-1}\asymp(\log t)^{-2}$ as
$t\to\infty$, establishing \eqref{eq:krein_asymp}.

Finally, since $f^*(x)=w(x_1)$ is independent of $(x_2,x_3)$, the
three-dimensional transition density of $Z_t$ factorizes into the
product of $p_{2t}^W$ and the standard two-dimensional Brownian heat
kernel $(8\pi t)^{-1}$ on $\R^2$ \cite{johnstonelalley1984}. The
expected total occupation time of a compact neighborhood of the origin
is governed by
\[
\int_2^\infty p_{2t}^W(0,0)\,(8\pi t)^{-1}\dd t
\asymp\int_2^\infty\frac{\dd t}{t(\log t)^2}<\infty.
\]
Hence $Z_t$ is transient on $\R^3$, and $\delta_F$ is inadmissible by
Theorem~\ref{thm:brown}(b).
\end{proof}

\begin{thebibliography}{99}

\bibitem{stein1956}
C.~Stein, \emph{Inadmissibility of the usual estimator for the mean of a
multivariate normal distribution}, in Proc.\ Third Berkeley Symp.\
Math.\ Statist.\ Probab., Vol.~1, Univ.\ California Press (1956),
pp.~197--206.

\bibitem{brown1971}
L.~D.~Brown, \emph{Admissible estimators, recurrent diffusions, and
insoluble boundary value problems}, Ann.\ Math.\ Statist.\ \textbf{42}
(1971), 855--903.

\bibitem{komaki2006}
F.~Komaki, \emph{Shrinkage priors for Bayesian prediction}, Ann.\
Statist.\ \textbf{34} (2006), 808--819.

\bibitem{komaki2015}
F.~Komaki, \emph{Asymptotic properties of Bayesian predictive densities
when the distributions of data and target variables are different},
Bayesian Anal.\ \textbf{10} (2015), 31--51.

\bibitem{amari2000}
S.~Amari and H.~Nagaoka, \emph{Methods of Information Geometry},
Translations of Mathematical Monographs, Vol.~191, American Mathematical
Society (2000).

\bibitem{bakryemery1985}
D.~Bakry and M.~\'Emery, \emph{Diffusions hypercontractives}, in
S\'eminaire de Probabilit\'es XIX, Lecture Notes in Math., Vol.~1123,
Springer (1985), pp.~177--206.

\bibitem{perelman2002}
G.~Perelman, \emph{The entropy formula for the Ricci flow and its
geometric applications}, arXiv:math/0211159 (2002).

\bibitem{grigoryan1999}
A.~Grigor'yan, \emph{Analytic and geometric background of recurrence and
non-explosion of the Brownian motion on Riemannian manifolds}, Bull.\
Amer.\ Math.\ Soc.\ (N.S.) \textbf{36} (1999), 135--249.

\bibitem{efron1975}
B.~Efron, \emph{Defining the curvature of a statistical problem (with
applications to second order efficiency)}, Ann.\ Statist.\ \textbf{3}
(1975), 1189--1242.

\bibitem{lyonssullivan1984}
T.~Lyons and D.~Sullivan, \emph{Function theory, random paths and
covering spaces}, J.\ Differential Geom.\ \textbf{19} (1984), 299--323.

\bibitem{grigoryan1985}
A.~A.~Grigor'yan, \emph{The existence of positive fundamental solutions
of the Laplace equation on Riemannian manifolds}, Mat.\ Sb.\ (N.S.)
\textbf{128(170)} (1985), 354--363; English transl.\ Math.\ USSR-Sb.\
\textbf{56} (1987), 349--358.

\bibitem{grigoryan1991}
A.~A.~Grigor'yan, \emph{The heat equation on noncompact Riemannian
manifolds}, Mat.\ Sb.\ \textbf{182} (1991), 55--87; English transl.\
Math.\ USSR-Sb.\ \textbf{72} (1992), 47--77.

\bibitem{saloffcoste1992}
L.~Saloff-Coste, \emph{A note on Poincar\'e, Sobolev, and Harnack
inequalities}, Internat.\ Math.\ Res.\ Notices \textbf{1992} (1992),
27--38.

\bibitem{sturm1996}
K.-T.~Sturm, \emph{Analysis on local Dirichlet spaces. III. The
parabolic Harnack inequality}, J.\ Math.\ Pures Appl.\ (9) \textbf{75}
(1996), 273--297.

\bibitem{kilpelainen1994}
T.~Kilpel\"ainen, \emph{Weighted Sobolev spaces and capacity}, Ann.\
Acad.\ Sci.\ Fenn.\ Ser.\ A I Math.\ \textbf{19} (1994), 95--113.

\bibitem{bjorn2011}
A.~Bj\"orn and J.~Bj\"orn, \emph{Nonlinear Potential Theory on Metric
Spaces}, EMS Tracts in Mathematics, Vol.~17, European Mathematical
Society (2011).

\bibitem{kasahara1975}
Y.~Kasahara, \emph{Spectral theory of generalized second order
differential operators and its applications to Markov processes},
Japan.\ J.\ Math.\ (N.S.) \textbf{1} (1975), 67--84.

\bibitem{kotani1982}
S.~Kotani and S.~Watanabe, \emph{Krein's spectral theory of strings and
generalized diffusion processes}, in Functional Analysis in Markov
Processes, Lecture Notes in Math., Vol.~923, Springer (1982),
pp.~235--259.

\bibitem{johnstonelalley1984}
I.~M.~Johnstone and S.~P.~Lalley, \emph{On independent statistical
decision problems and products of diffusions}, Z.\ Wahrsch.\ Verw.\
Gebiete \textbf{68} (1984), 29--47.

\end{thebibliography}

\end{document}